\documentclass[11pt,a4paper]{article}

% --- Mathematics (must precede unicode-math) ---
\usepackage{amsmath}
\usepackage{mathtools}

% --- Fonts (lualatex) ---
\usepackage{fontspec}
\usepackage{unicode-math}
\setmathfont{KpMath-Regular.otf}
\setmathfont{KpMath-Bold.otf}[version=bold]

% --- Microtypography (after all font setup) ---
\usepackage{microtype}

% --- Colors & page layout ---
\usepackage[dvipsnames]{xcolor}
\usepackage[margin=25mm, top=30mm, bottom=30mm]{geometry}

% --- Headers & footers ---
\usepackage{fancyhdr}
\pagestyle{fancy}
\fancyhf{}
\fancyhead[L]{\small\itshape Coupled Foldy-Lax T-Matrix}
\fancyhead[R]{\small\itshape Nestor (2026)}
\fancyfoot[C]{\thepage}
\renewcommand{\headrulewidth}{0.4pt}
\renewcommand{\footrulewidth}{0pt}

% --- Section numbering ---
\setcounter{secnumdepth}{3}

% --- Tables ---
\usepackage{booktabs}
\usepackage{array}

% --- Captions ---
\usepackage[font=small, labelfont=bf, labelsep=period]{caption}

% --- Spacing ---
\usepackage{setspace}
\setstretch{1.15}
\setlength{\parskip}{0.4em}

% --- Hyperlinks (last before macros) ---
\usepackage[colorlinks=true, linkcolor=blue!60!black,
            citecolor=blue!60!black, urlcolor=blue!60!black]{hyperref}

% --- Macros (reused from cubic_tmatrix.tex) ---
\newcommand{\bG}{\mathbf{G}}
\newcommand{\bI}{\mathbf{I}}
\newcommand{\bx}{\mathbf{x}}
\newcommand{\beR}{\hat{\mathbf{e}}_R}
\newcommand{\dd}{\mathrm{d}}
\newcommand{\order}[1]{\mathcal{O}\!\left(#1\right)}

\newcommand{\bT}{\mathbf{T}}
\newcommand{\bE}{\mathbf{E}}
\newcommand{\bS}{\mathbf{S}}
\newcommand{\Dlam}{\Delta\lambda}
\newcommand{\Dmu}{\Delta\mu}
\newcommand{\Drho}{\Delta\rho}
\newcommand{\PV}{\operatorname{P.V.}}
\newcommand{\re}{\operatorname{Re}}
\newcommand{\im}{\operatorname{Im}}

% --- New macros for this document ---
\newcommand{\bP}{\mathbf{P}}
\newcommand{\bC}{\mathbf{C}}
\newcommand{\bH}{\mathbf{H}}
\newcommand{\bM}{\mathbf{M}}
\newcommand{\bF}{\mathbf{F}}
\newcommand{\bpsi}{\boldsymbol{\psi}}
\newcommand{\beps}{\boldsymbol{\varepsilon}}

\begin{document}

% --- Title block ---
\thispagestyle{plain}
\begin{center}
\vspace*{1em}
{\LARGE\bfseries Coupled Foldy-Lax T-Matrix:\\[0.3em]
$(\Drho,\,\Dlam,\,\Dmu)$ Resonance Extension}\\[1.8em]
{\large T.\,M.\ Nestor}\\[0.5em]
{\itshape Research School of Earth Sciences, Australian National University}\\[2em]
\end{center}

% =============================================================================
\section{Introduction}
% =============================================================================

The Rayleigh-regime T-matrix for a cubic scatterer (see companion
document \emph{Cubic T-Matrix with the Full Elastodynamic Green's
Tensor}) captures density, bulk modulus, and shear modulus contrasts
through the self-consistent effective contrasts $\Drho^*$,
$\Dlam^*$, and $\Dmu^*$.  However, the displacement-only Foldy-Lax
formulation that extends this to the resonance regime ($ka \sim 1$)
captures only density scattering:\ the $3 \times 3$ sub-cell T-matrix
$\bT_\text{sub} = \omega^2 \Drho^* V_\text{sub}\,\bI_3$ has no
mechanism to transmit stiffness contrasts between sub-cells, because
stiffness scattering acts on the \emph{strain} field, not the
displacement field.

To include all three contrast types ($\Drho$, $\Dlam$, $\Dmu$) in
the internal Foldy-Lax system, the state vector at each sub-cell must
be enlarged from displacement alone (3~dof) to displacement plus Voigt
strain (9~dof).  The resulting $9N \times 9N$ coupled system---where
$N = n^3$ is the number of sub-cells---is the subject of this
document.

The coupled formulation:
\begin{enumerate}
  \item Subdivides the cube of half-width~$a$ into $n^3$ sub-cells,
    each of half-width $a_\text{sub} = a/n$, in the Rayleigh regime.
  \item Assigns a $9 \times 9$ local T-matrix to each sub-cell
    (density block $+$ stiffness block).
  \item Builds the $9 \times 9$ inter-sub-cell propagator from the
    Green's tensor $G_{ij}$, its first derivative $G_{ij,k}$, and its
    second derivative $G_{ij,kl}$.
  \item Solves the $9N \times 9N$ Foldy-Lax system and extracts the
    $9 \times 9$ composite T-matrix.
  \item Converts to the physical $6 \times 6$ displacement-traction
    T-matrix via the volume-integrated self-Green's tensor.
\end{enumerate}
The implementation is activated by passing \texttt{coupled=True} to
\texttt{compute\_resonance\_tmatrix()}.

% =============================================================================
\section{Green's Tensor and Derivatives}
\label{sec:greens}
% =============================================================================

\subsection{Kupradze representation}

The elastodynamic Green's tensor for a homogeneous isotropic full
space satisfies~\eqref{eq:pde_coupled} with the Kupradze solution
\begin{equation}
  G_{ij}(\bx)
  = P\bigl[\varphi(r)\,\delta_{ij}
    + \psi(r)\,\gamma_i\gamma_j\bigr],
  \qquad
  P = \frac{1}{4\pi\rho\omega^2},
  \label{eq:Kupradze}
\end{equation}
where $r = |\bx|$, $\gamma_i = x_i / r$, and
\begin{equation}
  \mu\,\nabla^2 G_{ij}
  + (\lambda + \mu)\,\partial_i\partial_k G_{kj}
  + \rho\omega^2 G_{ij}
  = -\delta_{ij}\,\delta(\bx).
  \label{eq:pde_coupled}
\end{equation}

The radial functions $\varphi$ and $\psi$ are sums of terms of the
form $c \cdot e^{ikr} \cdot r^p$ with $p \in \{-1,-2,-3\}$:
\begin{align}
  \varphi(r) &= k_S^2\,\frac{e^{ik_S r}}{r}
    + ik_S\,\frac{e^{ik_S r}}{r^2}
    - \frac{e^{ik_S r}}{r^3}
    - ik_P\,\frac{e^{ik_P r}}{r^2}
    + \frac{e^{ik_P r}}{r^3}, \label{eq:phi} \\[6pt]
  \psi(r) &= -k_S^2\,\frac{e^{ik_S r}}{r}
    - 3ik_S\,\frac{e^{ik_S r}}{r^2}
    + 3\,\frac{e^{ik_S r}}{r^3}
    + k_P^2\,\frac{e^{ik_P r}}{r}
    + 3ik_P\,\frac{e^{ik_P r}}{r^2}
    - 3\,\frac{e^{ik_P r}}{r^3}. \label{eq:psi}
\end{align}

\subsection{Derivative rules}

For each term $c \cdot e^{ikr} \cdot r^p$, the radial derivatives
follow from the chain rule on composite exponential-power
functions:
\begin{align}
  \frac{\dd}{\dd r}\bigl[e^{ikr} r^p\bigr]
  &= e^{ikr}\,r^{p-1}(ikr + p), \label{eq:d1} \\[4pt]
  \frac{\dd^2}{\dd r^2}\bigl[e^{ikr} r^p\bigr]
  &= e^{ikr}\,r^{p-2}\bigl[(ikr)^2 + 2p(ikr)
    + p(p-1)\bigr]. \label{eq:d2}
\end{align}
These are applied term-by-term to $\varphi$ and $\psi$ to obtain
$\varphi'$, $\psi'$, $\varphi''$, $\psi''$.

\subsection{First derivative \texorpdfstring{$G_{ij,k}$}{G\_ij,k}}

Using $\partial_k r = \gamma_k$ and
$\partial_k \gamma_i = (\delta_{ik} - \gamma_i\gamma_k)/r$:
\begin{equation}
  G_{ij,k} = P\Bigl[
    c_1\,\delta_{ij}\gamma_k
    + c_2\,\gamma_i\gamma_j\gamma_k
    + c_3\bigl(\delta_{ik}\gamma_j
      + \delta_{jk}\gamma_i\bigr)
  \Bigr],
  \label{eq:Gd}
\end{equation}
with three scalar coefficients:
\begin{equation}
  c_1 = \varphi', \qquad
  c_2 = \psi' - \frac{2\psi}{r}, \qquad
  c_3 = \frac{\psi}{r}.
  \label{eq:c123}
\end{equation}

\subsection{Second derivative \texorpdfstring{$G_{ij,kl}$}{G\_ij,kl}}

The second spatial derivative involves seven independent tensor
structures:
\begin{equation}
  \boxed{
  \begin{aligned}
    G_{ij,kl} = P\Bigl[
    &\;t_1\,\delta_{ij}\delta_{kl}
    + t_2\,\delta_{ij}\gamma_k\gamma_l
    + t_3\,\gamma_i\gamma_j\delta_{kl} \\[4pt]
    &+ t_4\bigl(\delta_{ik}\delta_{jl}
      + \delta_{jk}\delta_{il}\bigr)
    + t_3\bigl(\delta_{il}\gamma_j\gamma_k
      + \delta_{jl}\gamma_i\gamma_k \\
    &\quad + \delta_{ik}\gamma_j\gamma_l
      + \delta_{jk}\gamma_i\gamma_l\bigr)
    + t_7\,\gamma_i\gamma_j\gamma_k\gamma_l
    \Bigr],
  \end{aligned}
  }
  \label{eq:Gdd}
\end{equation}
with five independent scalar coefficients (note that $t_3$ appears
in three of the seven structures):

\begin{table}[ht]
\centering
\caption{Scalar coefficients for the seven tensor structures
  in $G_{ij,kl}$.}
\label{tab:Gdd_coeffs}
\begin{tabular}{@{}clll@{}}
\toprule
Coeff.\ & Expression & Tensor structures \\
\midrule
$t_1$ & $\varphi'/r$
  & $\delta_{ij}\delta_{kl}$ \\[4pt]
$t_2$ & $\varphi'' - \varphi'/r$
  & $\delta_{ij}\gamma_k\gamma_l$ \\[4pt]
$t_3$ & $\psi'/r - 2\psi/r^2$
  & $\gamma_i\gamma_j\delta_{kl}$, \;4 mixed terms \\[4pt]
$t_4$ & $\psi/r^2$
  & $\delta_{ik}\delta_{jl} + \delta_{jk}\delta_{il}$ \\[4pt]
$t_7$ & $\psi'' - 5\psi'/r + 8\psi/r^2$
  & $\gamma_i\gamma_j\gamma_k\gamma_l$ \\
\bottomrule
\end{tabular}
\end{table}

% =============================================================================
\section{Voigt Contraction}
\label{sec:voigt}
% =============================================================================

The coupled system requires contracting the Green's tensor derivatives
into the 6-component Voigt strain space.

\subsection{Voigt ordering}

The Voigt ordering maps symmetric pairs of Cartesian indices to a
single index $\alpha \in \{1,\ldots,6\}$:
\begin{equation}
  \alpha \;\longleftrightarrow\; (p,q):
  \quad
  1 = (0,0),\;
  2 = (1,1),\;
  3 = (2,2),\;
  4 = (1,2),\;
  5 = (0,2),\;
  6 = (0,1),
  \label{eq:voigt_map}
\end{equation}
corresponding to
$(\varepsilon_{zz},\,\varepsilon_{xx},\,\varepsilon_{yy},\,
  2\varepsilon_{xy},\,2\varepsilon_{zy},\,2\varepsilon_{zx})$
in the $(z,x,y)$ coordinate convention (axis labelling
$1 \to z$, $2 \to x$, $3 \to y$).

\subsection{Voigt propagator blocks}

From $G_{ij,k}$ and $G_{ij,kl}$, three Voigt-contracted matrices
are formed:
\begin{align}
  C_{i\alpha} &= \begin{cases}
    G_{ip,p} & \text{if } \alpha \leftrightarrow (p,p), \\
    G_{ip,q} + G_{iq,p} & \text{if } \alpha \leftrightarrow (p,q),
      \;p \neq q,
  \end{cases}
  \label{eq:Cblock} \\[8pt]
  H_{\alpha j} &= \begin{cases}
    G_{pj,p} & \text{if } \alpha \leftrightarrow (p,p), \\
    G_{pj,q} + G_{qj,p} & \text{if } \alpha \leftrightarrow (p,q),
      \;p \neq q,
  \end{cases}
  \label{eq:Hblock} \\[8pt]
  S_{\alpha\beta} &= \text{double Voigt contraction of } G_{ij,kl},
  \label{eq:Sblock}
\end{align}
with dimensions $\bC \in \mathbb{C}^{3 \times 6}$,
$\bH \in \mathbb{C}^{6 \times 3}$,
$\bS \in \mathbb{C}^{6 \times 6}$.

The block $\bC$ maps a Voigt stress source to a displacement
response; $\bH$ maps a force source to a Voigt strain response;
$\bS$ maps a Voigt stress source to a Voigt strain response.

\paragraph{Reciprocity.}
By the symmetry $G_{ij,k} = G_{ji,k}$ of the Green's tensor, the
relationship $\bH = \bC^\top$ holds.  This is verified numerically
in the implementation.

% =============================================================================
\section{The 9\texorpdfstring{$\times$}{x}9 Coupled System}
\label{sec:coupled}
% =============================================================================

\subsection{State vector}

At each sub-cell~$n$, the state vector combines displacement and
Voigt strain:
\begin{equation}
  \bpsi_n
  = \bigl(u_1,\,u_2,\,u_3,\,
    \varepsilon_{V1},\,\ldots,\,\varepsilon_{V6}\bigr)_n
  \in \mathbb{C}^9.
  \label{eq:state}
\end{equation}

\subsection{Local T-matrix (9\texorpdfstring{$\times$}{x}9)}

Each sub-cell has a block-diagonal local T-matrix:
\begin{equation}
  \boxed{
  \bT_\text{loc}
  = \begin{pmatrix}
      \omega^2 \Drho^* V_\text{sub}\,\bI_3 & \mathbf{0}_{3 \times 6} \\[4pt]
      \mathbf{0}_{6 \times 3} & V_\text{sub}\,\Delta\mathbf{c}^*_V
    \end{pmatrix},
  }
  \label{eq:Tloc}
\end{equation}
where $V_\text{sub} = (2a_\text{sub})^3$ is the sub-cell volume,
$\Drho^* = \Drho \cdot A_u$ is the self-consistent effective density
contrast, and $\Delta\mathbf{c}^*_V$ is the $6 \times 6$ effective
stiffness contrast in Voigt form:
\begin{equation}
  \Delta\mathbf{c}^*_V
  = \begin{pmatrix}
      \mathbf{D}^* & \mathbf{0}_{3 \times 3} \\
      \mathbf{0}_{3 \times 3} & \mathbf{S}^*
    \end{pmatrix},
  \qquad
  D^*_{ij} = \begin{cases}
    \Dlam^* + 2\Dmu^*_\text{diag} & i = j \\
    \Dlam^* & i \neq j
  \end{cases},
  \qquad
  S^*_{ij} = 2\Dmu^*_\text{off}\,\delta_{ij}.
  \label{eq:Dcstar}
\end{equation}

The density block produces a scattered force proportional to
displacement; the stiffness block produces a scattered stress
perturbation proportional to strain.  Both use the self-consistent
effective contrasts from the Rayleigh T-matrix computation
(\texttt{effective\_contrasts.py}).

\subsection{Inter-sub-cell propagator}

The $9 \times 9$ propagator between sub-cells $m$ and $n$ ($m \neq n$)
is:
\begin{equation}
  \boxed{
  \bP_{mn}
  = \begin{pmatrix}
      G_{ij} & C_{i\alpha} \\[4pt]
      H_{\alpha j} & S_{\alpha\beta}
    \end{pmatrix}
  \bigg|_{\bx = \bx_m - \bx_n},
  }
  \label{eq:Pmn}
\end{equation}
evaluated at the separation vector $\bx_m - \bx_n$.
The diagonal blocks ($m = n$) are zero---the self-interaction is
absorbed into $\bT_\text{loc}$ via the Rayleigh self-consistent
amplification factors.

\subsection{9\textit{N}\texorpdfstring{$\times$}{x}9\textit{N}
  Foldy-Lax equation}

Assembling the block-diagonal local T-matrix $\tilde{\bT}
= \bI_N \otimes \bT_\text{loc}$ and the off-diagonal propagator
$\tilde{\bP}$, the coupled Foldy-Lax system reads:
\begin{equation}
  \boxed{
  \bigl(\bI_{9N} - \tilde{\bP}\,\tilde{\bT}\bigr)\,
    \bpsi_\text{exc} = \bpsi_\text{inc},
  }
  \label{eq:FoldyLax}
\end{equation}
where $\bpsi_\text{inc}$ is the incident field and
$\bpsi_\text{exc}$ is the exciting (self-consistent) field at all
sub-cells.

\subsection{Incident field construction}

To extract the full $9 \times 9$ composite T-matrix, nine
independent incident patterns are needed:
\begin{itemize}
  \item \textbf{Columns 0--2} (displacement inputs): uniform unit
    displacement $u_i^\text{inc}(\bx_m) = \hat{e}_p$, zero incident
    strain $\beps_V^\text{inc} = \mathbf{0}$.
  \item \textbf{Columns 3--8} (strain inputs): linear displacement
    variation $u_i^\text{inc}(\bx_m) = \varepsilon^0_{ij}
    (x_{m,j} - \bar{x}_j)$ plus uniform Voigt strain
    $\beps_V^\text{inc} = \hat{e}_\alpha$.
\end{itemize}
A phase factor $e^{i\mathbf{k} \cdot \bx_m}$ is applied to all
patterns for resonance-regime accuracy.

\subsection{Composite T-matrix extraction}

The $9 \times 9$ composite T-matrix is obtained by summing the
scattered contributions from all sub-cells:
\begin{equation}
  \bT_\text{comp}[:,\,q]
  = \sum_{n=1}^{N} \bT_\text{loc} \cdot
    \bpsi_\text{exc}[9n{:}9(n{+}1),\,q],
  \label{eq:Tcomp}
\end{equation}
where $q \in \{1,\ldots,9\}$ indexes the nine independent incident
patterns.  The result $\bT_\text{comp}$ maps incident
$(u_i,\,\varepsilon_{V\alpha})$ to scattered
$(F_i,\,\Delta\sigma^*_{V\alpha})$---force monopole and stress dipole
amplitudes.

% =============================================================================
\section{Self-Interaction and Volume Integrals}
\label{sec:self}
% =============================================================================

The self-interaction (diagonal blocks) of the Foldy-Lax propagator
requires the volume integrals of $G_{ij}$, $G_{ij,k}$, and
$G_{ij,kl}$ over the cube.  These are already computed for the
Rayleigh T-matrix; here we collect the relevant results for the
output conversion stage.

\subsection{First-derivative integral}

By the symmetry of the cube ($[-a,a]^3$) and the odd parity of
$G_{ij,k}$ under inversion $\bx \to -\bx$:
\begin{equation}
  \int_{\text{cube}} G_{ij,k}\,\dd V = 0.
  \label{eq:Gd_vanish}
\end{equation}
This means the displacement--strain and strain--displacement
cross-coupling vanishes in the self-Green's tensor: $\bC_\text{self}
= \mathbf{0}$, $\bH_\text{self} = \mathbf{0}$.

\subsection{Second-derivative integral}

The volume integral of $G_{ij,kl}$ decomposes under cubic symmetry
exactly as the integral tensor from the companion document:
\begin{equation}
  \int_{\text{cube}} G_{ij,kl}\,\dd V
  = A\,\delta_{ij}\delta_{kl}
  + B\bigl(\delta_{ik}\delta_{jl} + \delta_{il}\delta_{jk}\bigr)
  + C\,E_{ijkl},
  \label{eq:Iself}
\end{equation}
where $E_{ijkl} = \delta_{ij}\delta_{kl}\delta_{ik}$ is the cubic
tensor (nonzero only when $i=j=k=l$).

Each coefficient receives a static (Eshelby) and smooth (radiation)
contribution:
\begin{equation}
  A = A_\text{stat} + A_\text{smooth}, \qquad
  B = B_\text{stat} + B_\text{smooth}, \qquad
  C = C_\text{stat} + C_\text{smooth}.
  \label{eq:ABC_split}
\end{equation}

\subsubsection*{Static Eshelby coefficients}

The static contributions arise from the principal-value volume
integral of the Kelvin Green's tensor, reduced to surface integrals
over the cube faces via the divergence theorem:
\begin{align}
  A_\text{stat} &=
  \frac{(\sqrt{3} - 2\pi)\,\alpha^2 - \sqrt{3}\,\beta^2}
       {6\pi\rho\alpha^2\beta^2},
  \label{eq:Astat} \\[6pt]
  B_\text{stat} &=
  \frac{\alpha^2 - \beta^2}
       {2\sqrt{3}\,\pi\rho\alpha^2\beta^2},
  \label{eq:Bstat} \\[6pt]
  C_\text{stat} &=
  \frac{(-5\sqrt{3} + 2\pi)(\alpha^2 - \beta^2)}
       {6\pi\rho\alpha^2\beta^2}.
  \label{eq:Cstat}
\end{align}
Note that $C_\text{stat}/B_\text{stat} = -5 + 2\pi/\sqrt{3}
\approx -1.372$ is a pure geometric constant (shape-dependent,
material-independent).  Both $B_\text{stat}$ and $C_\text{stat}$
are proportional to $(\alpha^2 - \beta^2)$ and vanish in the
mathematical limit $\alpha \to \beta$, where the T-matrix becomes
isotropic (sphere-like).  For physical materials $\alpha > \beta$
always holds (stability requires $K > 0$, so
$\alpha/\beta = \sqrt{(\lambda + 2\mu)/\mu} > \sqrt{4/3}$).

\subsubsection*{Smooth radiation corrections}

The smooth (radiation) contributions come from integrating the
regular part of the Green's tensor and its derivatives over the cube.
At leading order in $ka$:
\begin{align}
  A_\text{smooth} &\approx
    -\frac{2i\,a^3\omega^3}{3\pi\beta^5\rho}
    + \order{a^5\omega^5},
  \label{eq:Asmooth} \\[6pt]
  B_\text{smooth} &\approx
    \frac{i\,a^3\omega^3(\alpha^5 - \beta^5)}
         {3\pi\alpha^5\beta^5\rho}
    + \order{a^5\omega^5},
  \label{eq:Bsmooth} \\[6pt]
  C_\text{smooth} &\approx
    \frac{2i\,a^7\omega^7(\beta^9 - \alpha^9)}
         {4725\pi\alpha^9\beta^9\rho}
    + \order{a^9\omega^9}.
  \label{eq:Csmooth}
\end{align}
Since $C_\text{smooth}$ first appears at $\order{(ka)^7}$ while
$C_\text{stat}$ is frequency-independent, the cubic anisotropy is
dominated by the static Eshelby depolarization in the Rayleigh
regime.

\subsection{Voigt form of the self-integral}

Contracting~\eqref{eq:Iself} into Voigt indices gives the
$6 \times 6$ matrix $\bS_V$:
\begin{equation}
  \bS_V = \begin{pmatrix}
    \mathbf{D}_S & \mathbf{0} \\
    \mathbf{0} & \mathbf{S}_S
  \end{pmatrix},
  \label{eq:SV}
\end{equation}
with D-block entries
\begin{equation}
  (D_S)_{ii} = A + 2B + C, \qquad
  (D_S)_{ij} = B \quad (i \neq j),
  \label{eq:DS}
\end{equation}
and S-block entries
\begin{equation}
  (S_S)_{ii} = 2(A + B), \qquad
  (S_S)_{ij} = 0 \quad (i \neq j).
  \label{eq:SS}
\end{equation}

The block-diagonal structure of $\bS_V$ (vanishing D--S coupling)
is a consequence of the cubic symmetry of the cube and the
decomposition~\eqref{eq:Iself} involving only the three independent
tensors $\delta_{ij}\delta_{kl}$,
$\delta_{ik}\delta_{jl} + \delta_{il}\delta_{jk}$, and $E_{ijkl}$.

% =============================================================================
\section{Near-Field Exact 6\texorpdfstring{$\times$}{x}6 T-Matrix}
\label{sec:nearfield}
% =============================================================================

The composite T-matrix $\bT_\text{comp}$ operates in the 9-component
$(u_i, \varepsilon_{V\alpha})$ space, but the physical observable is
the $6 \times 6$ T-matrix in the displacement-traction basis
$(u_z, u_x, u_y, t_{zz}, t_{xz}, t_{yz})$.  Converting between
these requires careful treatment of the self-interaction Green's
tensor.

\subsection{Problem: far-field traction approximation}

For the displacement-only ($3 \times 3$) formulation, the scattered
traction is approximated using the far-field radiation pattern
($t_i \approx ik c_{ijkl} G_{kl} \cdots$).  This fails for densely
packed sub-cells where near-field interactions dominate, and it
completely misses the stiffness scattering contribution to the
traction.

\subsection{Three-stage pipeline}

The near-field exact $6 \times 6$ T-matrix is constructed via a
three-stage pipeline:
\begin{equation}
  \boxed{
  \bT_{6 \times 6}
  = \bG_\text{self} \cdot \bT_\text{comp} \cdot \bM_\text{in},
  }
  \label{eq:T6x6}
\end{equation}
where:

\paragraph{Stage 1: Input conversion $\bM_\text{in}$ (9$\times$6).}
Maps the incident displacement-traction vector
$(u_z, u_x, u_y, t_{zz}, t_{xz}, t_{yz})$ to the 9-component state
$(u_i, \varepsilon_{V\alpha})$:
\begin{equation}
  \bM_\text{in}
  = \begin{pmatrix}
      \bI_3 & \mathbf{0}_{3 \times 3} \\[4pt]
      \multicolumn{2}{c}{\bS(k_x, k_y)}
    \end{pmatrix},
  \label{eq:Min}
\end{equation}
where $\bS(k_x, k_y)$ is the $6 \times 6$ strain extraction matrix
that maps $(u_z, u_x, u_y, t_{zz}, t_{xz}, t_{yz})$ to Voigt strain
$(\varepsilon_{zz}, \varepsilon_{xx}, \varepsilon_{yy},
  2\varepsilon_{xy}, 2\varepsilon_{zy}, 2\varepsilon_{zx})$.  The
strain extraction follows from the constitutive relation on a
$z = \text{const}$ surface:
\begin{align}
  \varepsilon_{zz}  &= \frac{t_{zz} - \lambda\,ik_x u_x
    - \lambda\,ik_y u_y}{\lambda + 2\mu},
  &\varepsilon_{xx}  &= ik_x\,u_x, \notag \\
  \varepsilon_{yy}  &= ik_y\,u_y,
  &2\varepsilon_{xy} &= ik_y\,u_x + ik_x\,u_y, \notag \\
  2\varepsilon_{zy} &= t_{yz} / \mu,
  &2\varepsilon_{zx} &= t_{xz} / \mu.
  \label{eq:strain_extract}
\end{align}

\paragraph{Stage 2: Composite T-matrix $\bT_\text{comp}$
  (9$\times$9).}
The coupled Foldy-Lax output, mapping incident
$(u_i, \varepsilon_{V\alpha})$ to scattered
$(F_i, \Delta\sigma^*_{V\alpha})$.

\paragraph{Stage 3: Output conversion $\bG_\text{self}$
  (6$\times$9).}
Converts the scattered source amplitudes
$(F_i, \Delta\sigma^*_{V\alpha})$ to the scattered
displacement-traction $(u_i, t_i)^\text{scat}$ via the
volume-integrated self-Green's tensor:
\begin{equation}
  \bG_\text{self}
  = \begin{pmatrix}
      \Gamma_0\,\bI_3 & \mathbf{0}_{3 \times 6} \\[4pt]
      \mathbf{0}_{3 \times 3} & \bC_\text{trac} \cdot \bS_V
    \end{pmatrix},
  \label{eq:Gself}
\end{equation}
where $\Gamma_0 = \int_\text{cube} G_{ii}\,\dd V / 3$ is the scalar
volume integral of the Green's tensor, and $\bS_V$ is the Voigt form
of the self-integral~\eqref{eq:SV}.

The block-diagonal structure follows from the vanishing of the
cross-coupling integrals ($\bC_\text{self} = \bH_\text{self}
= \mathbf{0}$, \S\ref{sec:self}).

\subsection{Traction extraction matrix}

The $3 \times 6$ matrix $\bC_\text{trac}$ extracts the traction on a
$z = \text{const}$ surface from the Voigt stress:
\begin{equation}
  \bC_\text{trac}
  = \begin{pmatrix}
      \lambda + 2\mu & \lambda & \lambda & 0 & 0 & 0 \\
      0 & 0 & 0 & 0 & 0 & \mu \\
      0 & 0 & 0 & 0 & \mu & 0
    \end{pmatrix},
  \label{eq:Ctrac}
\end{equation}
mapping $(\sigma_{zz}, \sigma_{xx}, \sigma_{yy}, \sigma_{xy},
\sigma_{zy}, \sigma_{zx})_V$ to $(t_{zz}, t_{xz}, t_{yz})$.

\subsection{Complete formula}

Combining all three stages:
\begin{equation}
  \bT_{6 \times 6}(k_x, k_y)
  = \begin{pmatrix}
      \Gamma_0\,\bI_3 & \mathbf{0} \\
      \mathbf{0} & \bC_\text{trac}\,\bS_V
    \end{pmatrix}
  \cdot \bT_\text{comp}
  \cdot \begin{pmatrix}
      \bI_3 & \mathbf{0} \\
      \multicolumn{2}{c}{\bS(k_x, k_y)}
    \end{pmatrix}.
  \label{eq:T6x6_full}
\end{equation}
This is the near-field exact result, valid at all frequencies and for
all contrast types.  The dependence on horizontal wavenumbers
$(k_x, k_y)$ enters only through the strain extraction
matrix~$\bS$.

% =============================================================================
\section{Implementation Notes}
\label{sec:implementation}
% =============================================================================

\subsection{Function map}

Table~\ref{tab:functions} maps the key functions in the
implementation to the equations in this document.

\begin{table}[ht]
\centering
\caption{Function-to-equation reference for the coupled Foldy-Lax
  implementation.}
\label{tab:functions}
\begin{tabular}{@{}lll@{}}
\toprule
Function & Equation & Size \\
\midrule
\texttt{elastodynamic\_greens\_deriv}
  & \eqref{eq:Kupradze}, \eqref{eq:Gd}, \eqref{eq:Gdd}
  & $3^2, 3^3, 3^4$ \\[2pt]
\texttt{\_voigt\_contract}
  & \eqref{eq:Cblock}--\eqref{eq:Sblock}
  & $3{\times}6, 6{\times}3, 6{\times}6$ \\[2pt]
\texttt{\_propagator\_block\_9x9}
  & \eqref{eq:Pmn}
  & $9{\times}9$ \\[2pt]
\texttt{\_sub\_cell\_tmatrix\_9x9}
  & \eqref{eq:Tloc}
  & $9{\times}9$ \\[2pt]
\texttt{\_build\_incident\_field\_coupled}
  & \S\ref{sec:coupled}, \S4.5
  & $9N{\times}9$ \\[2pt]
\texttt{\_solve\_coupled}
  & \eqref{eq:FoldyLax}, \eqref{eq:Tcomp}
  & $9N{\times}9N$ \\[2pt]
\texttt{\_volume\_integral\_voigt}
  & \eqref{eq:SV}--\eqref{eq:SS}
  & $6{\times}6$ \\[2pt]
\texttt{\_traction\_from\_voigt\_strain}
  & \eqref{eq:Ctrac}
  & $3{\times}6$ \\[2pt]
\texttt{\_voigt\_tmatrix\_near\_field}
  & \eqref{eq:T6x6}
  & $6{\times}6$ \\[2pt]
\texttt{strain\_from\_displacement\_traction}
  & \eqref{eq:strain_extract}
  & $6{\times}6$ \\[2pt]
\texttt{effective\_stiffness\_voigt}
  & \eqref{eq:Dcstar}
  & $6{\times}6$ \\
\bottomrule
\end{tabular}
\end{table}

\subsection{Solver options}

Two solvers are available:
\begin{itemize}
  \item \textbf{Direct (LU):} \texttt{neumann\_order = 0} (default).
    Exact solution via $\texttt{np.linalg.solve}$.  Cost:
    $\order{(9N)^3}$.
  \item \textbf{Truncated Neumann series:}
    \texttt{neumann\_order\;>\;0}.
    Iterative:
    $\bpsi_\text{exc} = \sum_{k=0}^{K}
      (\tilde{\bP}\tilde{\bT})^k \bpsi_\text{inc}$.
    Cost per iteration: $\order{(9N)^2}$; converges when
    $\|\tilde{\bP}\tilde{\bT}\| < 1$.
\end{itemize}

\subsection{Computational cost}

The coupled system has size $9N \times 9N$ compared to
$3N \times 3N$ for the displacement-only formulation.  The direct
solve cost scales as $(9N)^3 / (3N)^3 = 27$ times larger.
Propagator assembly scales as $9^2 / 3^2 = 9$ times larger per
sub-cell pair, with the additional cost of computing $G_{ij,k}$ and
$G_{ij,kl}$ (five radial functions vs.\ two).

For moderate $N$ (e.g.\ $n = 3$, $N = 27$, system size $243$), the
direct solve remains tractable.  For larger $N$, the Neumann series
or iterative Krylov solvers are recommended.

\subsection{Backward compatibility}

Setting \texttt{coupled=False} (default) recovers the
displacement-only $3N \times 3N$ system.  When $\Dlam = \Dmu = 0$
(density-only contrast), the coupled system produces the same
$3 \times 3$ displacement T-matrix as the uncoupled system, since
the stiffness block of $\bT_\text{loc}$ vanishes and the strain
channels decouple.

\subsection{Validation}

Two key consistency checks are built into the test suite:
\begin{enumerate}
  \item \textbf{Rayleigh limit ($n = 1$, coupled):}  With a single
    sub-cell ($N = 1$), no inter-sub-cell coupling exists and the
    composite T-matrix must equal the single-cell Rayleigh result.
  \item \textbf{Density-only limit ($\Dlam = \Dmu = 0$, coupled):}
    With zero stiffness contrast, the coupled $9 \times 9$ system
    must recover the displacement-only $3 \times 3$ T-matrix
    identically.
\end{enumerate}

\end{document}
